\labsection{6}{NLP Foundations: Handwriting OCR Evaluation}{sec:lab06-nlp-a}

\begin{labproject}{Evaluate OCR on real handwritten medicine names}
\textbf{Scenario.} A clinic wants to digitise handwritten medicine names before storage or search.\\
\textbf{Goal.} Build a reproducible OCR benchmark on 89 labelled photographs, test one controlled preprocessing change, and analyse the failures.\\
\textbf{Data.} \texttt{all\_datasets/lab06-ocr-handwritten-prescription.zip} --- one label CSV and 89 images. The notebook extracts it next to the archive.\\
\textbf{Source files.} \texttt{all\_experiments/lab06\_handwriting\_ocr\_worked.ipynb}.\\
\end{labproject}

\begin{labmode}
\texttt{pytesseract} is only a wrapper; the Tesseract \emph{engine} must be installed separately and the kernel restarted afterwards. The setup cell looks for it on \texttt{PATH} and in the default Windows install folder.
\begin{lstlisting}[style=aiterminal]
winget install --id UB-Mannheim.TesseractOCR    # Windows
sudo apt-get install tesseract-ocr              # Ubuntu / Debian
brew install tesseract                          # macOS
\end{lstlisting}
\end{labmode}

\noindent\textbf{Setup cell.}
\begin{lstlisting}[style=aipython]
import re
import shutil
import zipfile
from pathlib import Path

import numpy as np
import pandas as pd
import matplotlib.pyplot as plt
import pytesseract
from IPython.display import display
from PIL import Image, ImageOps

ROOT = Path.cwd().parent
DATA = ROOT / "all_datasets"
RESULTS = ROOT / "results" / "lab06"
RESULTS.mkdir(parents=True, exist_ok=True)

# The engine is found on PATH or in the default Windows install folder.
tesseract = shutil.which("tesseract") or r"C:\Program Files\Tesseract-OCR\tesseract.exe"
pytesseract.pytesseract.tesseract_cmd = tesseract
print("Tesseract", pytesseract.get_tesseract_version(), "at", tesseract)
\end{lstlisting}

\begin{labquestion}{1}{Validate the archive}
Extract the archive next to itself, load the label file, check that all 89 labelled images exist exactly once, report the number of unique labels, and show a 3-by-3 sample grid.
\end{labquestion}

\begin{labanswer}
\begin{lstlisting}[style=aipython]
IMAGES = DATA / "lab06-ocr-handwritten-prescription"
if not IMAGES.exists():
    zipfile.ZipFile(DATA / "lab06-ocr-handwritten-prescription.zip").extractall(IMAGES)

labels = pd.read_csv(IMAGES / "doctor_handwriting_labels.csv")
labels["path"] = [IMAGES / "img" / "img" / name for name in labels.filename]
assert len(labels) == 89 and labels.filename.is_unique and all(p.exists() for p in labels.path)
print("Labelled images:", len(labels), "| unique labels:", labels.label.nunique())

fig, axes = plt.subplots(3, 3, figsize=(10, 6))
for ax, row in zip(axes.flat, labels.head(9).itertuples()):
    ax.imshow(Image.open(row.path))
    ax.set_title(row.label, fontsize=9)
    ax.axis("off")
plt.tight_layout()
plt.savefig(RESULTS / "sample_grid.png", dpi=150)
\end{lstlisting}
\expectedoutput
\begin{lstlisting}[style=aiterminal]
Labelled images: 89 | unique labels: 61
\end{lstlisting}
One label file and 89 photographs, each present exactly once; the 89 labels hold 61 distinct medicine names, so some names appear several times in different handwriting. The grid shows what the engine faces: cursive, varied stroke width, uneven lighting and lined paper.
\end{labanswer}

\begin{labquestion}{2}{Run the baseline OCR}
Use one fixed Tesseract configuration for every raw image (LSTM engine, page-segmentation mode~7 = one text line) and normalise both prediction and label to lowercase letters and digits, so punctuation and spacing do not count as errors.
\end{labquestion}

\begin{labanswer}
\begin{lstlisting}[style=aipython]
OCR_CONFIG = "--oem 1 --psm 7"


def normalise(text):
    return re.sub(r"[^a-z0-9]", "", str(text).lower())


def run_ocr(paths, pipeline):
    """One row per image: the raw prediction and its normalised form."""
    rows = []
    for row, path in zip(labels.itertuples(), paths):
        prediction = pytesseract.image_to_string(Image.open(path), config=OCR_CONFIG).strip()
        rows.append({"filename": row.filename, "pipeline": pipeline, "truth": row.label, "prediction": prediction,
                     "truth_norm": normalise(row.label), "prediction_norm": normalise(prediction)})
    return pd.DataFrame(rows)


baseline = run_ocr(labels.path, "baseline")
display(baseline[["truth", "prediction", "prediction_norm"]].head(8))
\end{lstlisting}
\expectedoutput
\begin{lstlisting}[style=aiterminal]
          truth      prediction prediction_norm
0  cefiget 40mg
1    tab azimax             Ago             ago
2  tab toniflex  p53 Towid host    p53towidhost
3          nims            Nidu            nidu
4      dicloran        divlarew        divlarew
...
\end{lstlisting}
The first rows already show the two failure modes the lab will quantify: images where Tesseract returns nothing, and images where it returns text that shares only a few characters with the label. The configuration is now fixed, so any later improvement must come from the image, not from tuning the engine.
\end{labanswer}

\begin{labquestion}{3}{Measure recognition quality}
Implement Levenshtein edit distance and report, for the baseline: exact-match accuracy, mean normalised edit distance (distance divided by the longer string, so it lies in 0--1), character error rate (total edits divided by total reference characters; it can exceed~1), and the number of blank predictions.
\end{labquestion}

\begin{labanswer}
\begin{lstlisting}[style=aipython]
def levenshtein(a, b):
    """Minimum number of single-character edits that turn a into b."""
    previous = list(range(len(b) + 1))
    for i, char_a in enumerate(a, start=1):
        current = [i]
        for j, char_b in enumerate(b, start=1):
            current.append(min(current[j - 1] + 1, previous[j] + 1, previous[j - 1] + (char_a != char_b)))
        previous = current
    return previous[-1]


def add_distances(frame):
    frame = frame.copy()
    frame["exact"] = frame.truth_norm == frame.prediction_norm
    frame["edits"] = [levenshtein(t, p) for t, p in zip(frame.truth_norm, frame.prediction_norm)]
    longest = np.maximum(frame.truth_norm.str.len(), frame.prediction_norm.str.len()).clip(lower=1)
    frame["normalised_edit_distance"] = frame.edits / longest
    return frame


def summarise(frame):
    return {"pipeline": frame.pipeline.iat[0],
            "exact_match_accuracy": frame.exact.mean(),
            "mean_normalised_edit_distance": frame.normalised_edit_distance.mean(),
            "character_error_rate": frame.edits.sum() / frame.truth_norm.str.len().sum(),
            "blank_predictions": int((frame.prediction_norm == "").sum())}


baseline = add_distances(baseline)
display(pd.DataFrame([summarise(baseline)]).round(4))
\end{lstlisting}
\expectedoutput
\begin{lstlisting}[style=aiterminal]
   pipeline  exact_match_accuracy  mean_normalised_edit_distance  character_error_rate  blank_predictions
0  baseline                   0.0                         0.8559                0.8753                 39
\end{lstlisting}
Generic OCR on raw handwriting is poor: none of the 89 names is read exactly, the character error rate is close to~0.9, and 39 images yield no text at all. The metrics are reproducible because the normalisation and the distance are written out rather than taken from a library with hidden options.
\end{labanswer}

\begin{labquestion}{4}{Build one preprocessing pipeline}
Convert each image to grayscale, add a 20-pixel white border, cap the width at 1,200 pixels, and enlarge by 1.5$\times$ with bicubic resampling. Save the results under \texttt{results/lab06/preprocessed/}. The OCR configuration does not change.
\end{labquestion}

\begin{labanswer}
\begin{lstlisting}[style=aipython]
def preprocess(path, output):
    image = ImageOps.exif_transpose(Image.open(path)).convert("L")
    image = ImageOps.expand(image, border=20, fill=255)
    if image.width > 1200:
        image = image.resize((1200, round(image.height * 1200 / image.width)), Image.Resampling.LANCZOS)
    image = image.resize((round(image.width * 1.5), round(image.height * 1.5)), Image.Resampling.BICUBIC)
    image.save(output)
    return output


(RESULTS / "preprocessed").mkdir(exist_ok=True)
preprocessed_paths = [preprocess(path, RESULTS / "preprocessed" / f"{i:03d}.png") for i, path in enumerate(labels.path)]
example = Image.open(preprocessed_paths[0])
print("Prepared images:", len(preprocessed_paths), "| example size:", example.size, "mode:", example.mode)
\end{lstlisting}
\expectedoutput
\begin{lstlisting}[style=aiterminal]
Prepared images: 89 | example size: (1800, 638) mode: L
\end{lstlisting}
Each step targets a known weakness of Tesseract on photographs: grayscale removes colour noise, the border stops text touching the edge (which breaks line segmentation), the width cap bounds runtime, and the enlargement brings stroke height into the range the engine was trained on. Only the image changes; the engine settings do not.
\end{labanswer}

\begin{labquestion}{5}{Compare the two pipelines fairly}
Run OCR on the preprocessed images with the identical configuration, keep exactly the same 89 samples, save one row per image and pipeline to \texttt{predictions.csv}, and save both sets of metrics to \texttt{metrics.csv}.
\end{labquestion}

\begin{labanswer}
\begin{lstlisting}[style=aipython]
preprocessed = add_distances(run_ocr(preprocessed_paths, "grayscale_border_resize"))
pd.concat([baseline, preprocessed]).to_csv(RESULTS / "predictions.csv", index=False)

metrics = pd.DataFrame([summarise(baseline), summarise(preprocessed)]).set_index("pipeline")
metrics.to_csv(RESULTS / "metrics.csv")
display(metrics.round(4))
display((metrics.loc["grayscale_border_resize"] - metrics.loc["baseline"]).rename("change").round(3).to_frame())
\end{lstlisting}
\expectedoutput
\begin{lstlisting}[style=aiterminal]
                         exact_match_accuracy  mean_normalised_edit_distance  character_error_rate  blank_predictions
pipeline
baseline                               0.0000                         0.8559                0.8753                 39
grayscale_border_resize                0.0112                         0.7673                0.7883                 15
                               change
exact_match_accuracy            0.011
mean_normalised_edit_distance  -0.089
character_error_rate           -0.087
blank_predictions             -24.000
\end{lstlisting}
Preprocessing cuts the blank predictions from 39 to 15 and lowers the character error rate by about 0.09, but exact-match accuracy barely moves (one correct name out of 89): the engine now \emph{sees} text on most images yet still cannot read cursive. Because both pipelines were scored on the same 89 images with the same configuration, the difference is attributable to the preprocessing alone. Exact numbers depend on the Tesseract version installed; the direction of the changes does not.
\end{labanswer}

\begin{labquestion}{6}{Analyse the failures}
Take the worst preprocessed predictions by normalised edit distance, tag each one as \emph{no text detected} or \emph{character confusion}, save the table as \texttt{error\_analysis.csv}, and save a figure of the six hardest images.
\end{labquestion}

\begin{labanswer}
\begin{lstlisting}[style=aipython]
errors = preprocessed[~preprocessed.exact].sort_values(["normalised_edit_distance", "filename"], ascending=[False, True]).head(8)
errors = errors.assign(failure_type=np.where(errors.prediction_norm == "", "no text detected", "character confusion"))
errors[["filename", "truth", "prediction", "failure_type"]].to_csv(RESULTS / "error_analysis.csv", index=False)
display(errors[["truth", "prediction", "failure_type"]])

fig, axes = plt.subplots(2, 3, figsize=(11, 5.5))
for ax, row in zip(axes.flat, errors.head(6).itertuples()):
    ax.imshow(Image.open(IMAGES / "img" / "img" / row.filename))
    ax.set_title(f"Reference: {row.truth}\nOCR: {row.prediction[:28]}", fontsize=8)
    ax.axis("off")
plt.tight_layout()
plt.savefig(RESULTS / "ocr_failures.png", dpi=150)
print("Saved:", sorted(p.name for p in RESULTS.iterdir() if p.is_file()))
\end{lstlisting}
\expectedoutput
\begin{lstlisting}[style=aiterminal]
               truth prediction      failure_type
60            azitma             no text detected
61          tonoflox             no text detected
0       cefiget 40mg             no text detected
...
Saved: ['error_analysis.csv', 'metrics.csv', 'ocr_failures.png', 'predictions.csv', 'sample_grid.png']
\end{lstlisting}
The hardest cases split into two kinds. \emph{No text detected} is a segmentation failure --- faint strokes, low contrast or slanted writing that the line finder rejects. \emph{Character confusion} is a recognition failure --- a line is found but cursive shapes are mapped to the wrong printed letters. Preprocessing helps the first kind and barely touches the second, which is why exact-match accuracy stays low.

This is a benchmark of model--domain fit, not evidence that generic OCR suits clinical transcription. A medical system would need a handwriting-specific recogniser, representative data, external validation, privacy controls and mandatory human verification, and must never decide medication or dosage from uncertain text.
\end{labanswer}
